% !TeX root = ../papers/nc-p3.tex
\section{階数四の対称核と二項切断の全射性}
\label{sec:smooth-kernel-binary-transfer}
四辺が有限な標準有限Frobenius窓で、一方の対向位置が対称、反対階数9とする。
対称二乗の核の生成元が階数4なら、対称側の点座標を $E=A\otimes B$、$\dim A=\dim B=2$ とし、
核を二行二列の行列式に対応するテンソルに取れる。
従って四テンソルは二つの $A$ 位置の交換と二つの $B$ 位置の交換で別々に不変である。
これは $\operatorname{Sym}^2(A\otimes B)^*=
(\operatorname{Sym}^2 A^*\otimes\operatorname{Sym}^2 B^*)\oplus
(\wedge^2 A^*\otimes\wedge^2 B^*)$ の分解による。
反対側の点座標を $F,G$ とし、隣接分解を $\beta:E\otimes F\to H$、$\gamma:E\otimes G\to H^*$、$\dim H=6$ とする。
\begin{theorem}\label{thm:smooth-kernel-binary-transfer}
いずれかの二項因子、例えば $A$ において、ある $a\in A$ で
$\beta(a,-,-):B\otimes F\to H$ と $\gamma(a,-,-):B\otimes G\to H^*$ がともに全射なら、
窓は両対称反転窓である。長さ1の局所点の存在は不要である。
従ってこの有限開条件では全次数の正則性が既出の有限判定で決まる。
\end{theorem}
\begin{proof}[一つの作用素]
$a$ を第一基底に取り、二写像の二係数を $P_0,P_1$、$Q_0,Q_1$ とする。
$A$ の対称性から $P_0^t Q_1=P_1^t Q_0$。$Q_0$ の全射性により $P_1$ は
$P_0$ を経由し、両全射性から $P_1=T P_0,\ Q_1=T^t Q_0$ となる単一の $T\in\operatorname{End} H$ がある。
$T$ の各左固有汎関数は $R_\beta$ の元を与え、その偶位置の支持は固有値に対応する
$A^*$ の一次式と $B^*$ の積の二平面に入る。純関係がないので階数は常に2である。
同じ固有値の独立な二汎関数があれば、これらの右支持の和は $F^*$ 全体である。
実際、和が二次元以下なら二行二列の鉛筆に階数1がある。
三次元なら二関係の四つの行の依存は一つで、関係の係数と二行の間の二行列として非退化である
（階数1の依存なら非零の関係が階数1以下となる）。
基底を合わせると二関係は $u_1f_1+u_2f_2,\ u_1f_2+u_2f_3$ と書ける。
これらを順に $u_1,u_2$ と右からテンソルして加えたものは、両端対称の非零三次重なりで両端支持が2以下である。
$R_\beta,R_\beta^{\rm op}$ は連続し、全ての三次重なりは四テンソルの一文字縮約なので、
その点での辺の評価階数が2以下となり有限性に反する。$\gamma$ についても同じである。
固有空間の次元が3以上なら、対応する $a'\in\mathbb P A$ の $P(a')$ は階数3以下となり、
$\{a'\}\times\mathbb P B\times\mathbb P F$ の四次元上で三つ以下の双一次式を課すだけで
正次元の辺の点対ができる。従って各固有値のJordanブロックは高々二つである。
\end{proof}

\begin{proof}[巡回性と二次行列]
不変因子を $p$ および $q\mid p$ とし、巡回の場合 $q=1$ とする。
$m=\deg p,\ e=\deg q,\ m+e=6$ で、$H=D_p\oplus D_q$、$D_h=k[t]/(h)$ と書く。
$T$ は $t$ 倍である。各モニック $h$ に対し剰余代表の最高次係数（次数 $\deg h-1$）を取る
$\ell_h$ で積を対合させ、双対も同じ表示に取る。
$P_0,Q_0$ の $D_h$ 成分を $B$ の二基底で $U_{h,0},U_{h,1}$ と $V_{h,0},V_{h,1}$ と書く。
前者は $F^*$、後者は $G^*$ に値を持つ $D_h$ ベクトルである。$W_h=U_{h,0}V_{h,1}^t-U_{h,1}V_{h,0}^t$ とする。
第二の二項因子の対称性から $W=W_p+(p/q)W_q$ の各成分は $D_p$ 上で $1,t,t^2$ に直交する
（$q=1$ なら第二項なし）。実際 $\ell_q(s)=\ell_p((p/q)s)$ であり、二つの $A$ 変数の積の係数を読めばよい。
$m\ge3$ なので直交部分は次数 $m-4$ 以下の通常の多項式全体である。次数上限が負なら零と読む。
以下この代表 $W(t)$ を使う。

$e>0$ なら $m\le5$ で、$W$ の三次小行列式は $p/q$ で割れるが次数は $3(m-4)<m-e$ である
（$m=3$ では $W=0$）。従って全て恒等零となる。
一方、$q$ の根で $p,q$ の重複度を $l_1\ge l_2>0$ とする。
その局所環で両不変因子の成分を規定次数まで持ち上げると、先の固有空間の議論により
四列 $U_{p,0},U_{p,1},U_{q,0},U_{q,1}$ は剰余体で独立であり、$V$ の四列も独立である。
従って $W\bmod (t-\lambda)^{l_1}$ は左右可逆変換で、二次交代ブロックと
$(t-\lambda)^{l_1-l_2}$ 倍の二次交代ブロック（単元倍を許す）の直和である。
ある三次小行列式はこの局所環で非零であり矛盾する。
従って $T$ は巡回であり、$m=6$、$W$ は二次以下になる。

全三次小行列式は六次までで $p$ で割れるので $\operatorname{adj} W=p C$、$C$ は定数行列である。
$W$ が一般に可逆なら $C$ は可逆で $W$ は定数可逆行列のスカラー多項式倍となるが、
各根での階数は2である。一般階数3なら $C$ の非零行・列から定数左右核が生じる。
例えば $b^tW=0$ なら、全ての局所環で $V_{p,0},V_{p,1}$ は独立な直和因子を張るから
$b^t U_{p,0}=b^t U_{p,1}=0$ が $D_p$ で成り立つ。この文字は $P_0$ と $P_1=TP_0$ の双方で消え、
一因子全支持に反する。一般階数は従って2である。この議論は $p$ が重根だけの場合にも成立する。
\end{proof}

\begin{proof}[二つの四文字の同定]
$W(t)$ を二次に斉次化し、両側の一般像を射影直線上で飽和させた階数2の部分束を
$S_F\subset F^*\otimes\mathcal O$, $S_G\subset G^*\otimes\mathcal O$ とする。
$S_F$ の双対の次数を $d_F$ とすると、大域生成と直線上の分裂から $h^0(S_F^*)=2+d_F$。
先の定数核の排除により $F\to H^0(S_F^*)$ は単射で、$d_F\ge2$。
同じく $d_G\ge2$ である。一方、$W$ は $S_G^*\to S_F(2)$ を誘導し、行列式は
次数 $4-d_F-d_G$ の非零切断である。従って $d_F=d_G=2$、誘導写像は全点で同型、
二つの大域切断の単射も同型である。

$\det S_F\simeq\mathcal O(-2)$ を選ぶ。$W$ による同型と階数2の外積対合により
$S_G\simeq S_F$ と同定して $W$ を $S_F$ 上の非退化交代二テンソル（捻り2）にできる。
両者の双対の全切断が四文字そのものなので、この束同型は周囲の定数線型同型を誘導し、
両埋込みも双対評価で対応する。従って定数同定後の $W(t)$ は全域で交代である。
$D_p$ に戻ると、二列ずつの $U,V$ は同じ階数2直和因子を張る。
局所的に $V=U L$ と書けば $U\bigl(\begin{smallmatrix}0&1\\-1&0\end{smallmatrix}\bigr)L^t U^t$ の交代性から
$L$ は単元スカラー行列である。全局所成分を合わせ $V=h U$、$h\in D_p^\times$ と書ける。
従って $\gamma$ は相手四文字を同定した $\beta$ に対し、完全対称積対合
$(s,t)\mapsto\ell_p(hst)$ で与えられる。これは隣接二組の同時交換で四テンソルを保ち、
既知の偶位置の交換と合わせて奇位置の交換も保つ。これで両対称性が従う。
\end{proof}

\begin{theorem}\label{thm:smooth-symmetric-kernel-all}
同じ有限点グラフで対称核の生成元が階数4なら、どちらかの二項因子が前定理の同時全射条件を満たすか、
元代数が指数成長するかのいずれかである。従ってこの滑らかな対称核の全域で正則性は有限に決まる。
特に長さ1の局所成分を持つ反対階数9では、片側対称表示で残る未判定の二次形式は階数3だけである。
\end{theorem}
\begin{proof}[一つの写像の二方向の階数]
まず $A$ を動かす $P(a):B\otimes F\to H$ の一般階数は5以上である。
4以下なら各 $a$ で $\mathbb P B\times\mathbb P F$ の四次元に高々四つの超平面を課すため、
各 $a$ に辺の点対が存在して有限性に反する。$B$ 方向でも、また $\gamma$ でも同様である。

$\beta$ は両方向とも一般階数5にはなれないことを証明する。
そうと仮定し、左核の原始斉次ベクトルをそれぞれ $l(a),n(b)\in H^*$、次数を $k,h$ とする。
五次小行列式で表示できるので $k,h\le5$、原始性により各射影直線で基点はない。
一般の $(a,b)$ で $\beta(a,b,F)$ の階数は4で、その零化空間は $l(a),n(b)$ の二元で張られる。
実際、後二者が一般に比例すれば両方が一定となり $\beta$ の全射性に反する。
$6\times4$ 行列 $\beta(a,b,-)$ の四次小行列式は二次数 $(4,4)$ で最大公約因子はない（退化点は有限）。
従ってこれらの相補Plücker座標は $l(a)\wedge n(b)$ をその成分の最大公約因子で割ったものに比例する。
特に $k,h\ge4$ である。

$P_0,P_1$ の共通核を割った入力次元を $r$ とする。$r\ge6$ なら一般に階数5である六列を独立な入力から取る。
$6\times6$ 鉛筆の随伴行列の非零行は $l(a)$ の多項式倍であり、
次数が5なので、この六列の非零右核ベクトルが次数 $5-k\le1$ で得られる
（随伴行列を共通の左核でくくる）。定数核は既に割ったため、一次の
$a_0u_0+a_1u_1$ である。$P_0u_0=P_1u_1=0$, $v=P_0u_1=-P_1u_0\ne0$ となる。
ところが $P_0^t Q_1=P_1^tQ_0$ にこの二入力を入れると $v$ は $Q_0,Q_1$ の両像を消す。
$\gamma$ の全射性に反するので $r=5$、他方向でも同じである。

従って $B\otimes F$ の三次元共通核 $K$ には階数1元がない。あれば $a$ 全体で辺の零点となる。
この核は基底変更で $B\otimes S^3 B\to S^4 B$ の乗法核である。
実際、$\mathbb P(B^*)$ 上で評価した $K\otimes\mathcal O(-1)\to F\otimes\mathcal O$ は
全点で階数3、余核は $\mathcal O(3)$ で、大域切断で $F$ はその全切断に等しい。
完全三次系の線形関係は二変数の通常の乗法核である（体積形式で $B$ とその双対を同定してよい）。
よって $\beta(a,b,f)=\widetilde P(a)(bf)$、$\widetilde P$ は五列の鉛筆である。他方向にも同じ記述がある。

滑らかな二次曲面 $\mathbb P A\times\mathbb P B$ 上の点グラフの制限長は $20-4k$ である。
実際、$l$ の核束を $K_l=\ker(H\otimes\mathcal O\to\mathcal O(k))$ とし、
$\mathbb P A\times\mathbb P B\times\mathbb P F$ に引き戻すと、辺の零点は
階数5束 $K_l(1,1,1)$ の切断の零点である。有限なので空の場合も含めて最高Chern数が長さとなる。
各超平面類を $x,y,z$ とすると $(x+y+z)^5-kx(x+y+z)^4$ の $xyz^3$ 係数は $20-4k$ である。
射影はグラフ上で閉埋込みなので他方向の計算と比較して $k=h$ となる。
$k=4$ なら $\widetilde P(a)$ の最大小行列式は $l(a)$ の各成分と同じ一次式の積である。
零点上では $bf$ の単元成分があるため、$\widetilde P(a)(bf)=0$ により全最大小行列式が
スキーム上で消える。共通の一次式が消え、グラフの制限はこの一つの被約な母線に入る。
逆方向の一母線にも入るため長さは高々1となり、計算値4に反する。

$k=5$ なら $\widetilde P$ の階数は全点で5、束の余核は
$5\mathcal O(-1)\to6\mathcal O\to\mathcal O(5)$ である。
従って $l$ は完全五次系による有理正規曲線であり、$n$ も同様である。
$l\wedge n$ の最大公約因子は二次数 $(1,1)$ なので両曲線は一致する
（一方の変数を固定した成分は、他側が埋込みなのでありえない）。
二つのパラメータを射影同型で同定し $H=S^5 A$ と書くと、$\beta(a,b,f)$ は
全ての $a,b$ についてそれぞれが指定する二つの一次因子を持つ五次式である。
二つの普遍一次因子は独立な変数を用い互いに素なのでその積で割れる。
$\beta$ 自身は $a,b$ にそれぞれ一次で、商は $f$ のみに依存する。
従って偶位置の $A\otimes B$ は対称二乗を経由し、一文字全支持に反する。
これで $\beta$ の両方向同時の階数5を除外した。$\gamma$ にも同じ証明を適用できる。
\end{proof}

\begin{proof}[交叉する二方向の非全射]
同時全射の方向が一つもなければ、入れ替えを除いて $A$ 方向で $\beta$ の一般階数が5、
$\gamma$ が6、$B$ 方向では逆となる。一般の $a$ で $P(a)v=0$ とすると、
対称性と $Q(a)$ の全射により全 $a'$ で $P(a')v=0$ である。従って共通核の次元は3で、
先ほどの乗法核の同定により $F=S^3B$、$\beta$ は $B\otimes S^3B\to S^4B$ を経由する。
反対側では $G=S^3A$ で、$\gamma$ は $A\otimes S^3A\to S^4A$ を経由する。
隣接関係の四列は $R_\beta,R_\beta^{\rm op},R_\gamma,R_\gamma^{\rm op}$ である。

経路空間の双対を各点文字のテンソルに取り、二項多項式で表す。
各 $E,F$ 辺では、偶文字の $B$ 変数と奇文字の $B$ 変数の角括弧があれば二変数を掛ける写像で零となる。
$G,E$ 辺では $A$ 変数どうしの角括弧で同じことが成り立つ。
全辺の括弧を掛けると、内部の偶文字では $A,B$ の各一次をちょうど使い、内部の奇文字では
三次のうち二次しか使わない。各内部奇位置に任意の一次形式を掛け、端の不足次数を固定した非零形式で埋める。
得たテンソルは全隣接関係を消す。多項式環で固定した非零積の乗法は単射なので、内部奇位置の数を $s$ として
同じ長さに $2^s$ 個の線型独立元がある。これは標準Hilbert増大を排除する。
\end{proof}
